%% ch08 --- Droga do chaosu i liczba, której nikt się nie spodziewał
%%
%% Napisany we wrześniu 2026. Każda liczba, której czytelnik nie policzy w
%% pamięci, pochodzi z code/measure/ch08.py, który importuje TE SAME funkcje,
%% które drukują listingi (code/ch08/), więc strona i listingi nie mogą się
%% rozejść. Rozwidlenia i punkty orientacyjne odwzorowania logistycznego
%% liczone są tylko przez + - * i /, więc ich cyfry są takie same na każdej
%% maszynie; odwzorowanie sinusowe woła math.sin, więc to, co strona z niego
%% drukuje, jest zaokrąglone do dwóch miejsc (czterech dla miejsca, gdzie
%% kończą się jego podwojenia) i sprawdzone, że nie leży na granicy
%% zaokrąglenia.
%%
%% CZEGO TEN ROZDZIAŁ NIE MOŻE POWIEDZIEĆ. Cztery zachowania odwzorowania
%% logistycznego i utrata stabilności punktu stałego przy r = 3 należą do
%% rozdziału 5; tu są przypomniane, nie uczone. Wykładnik Lapunowa w funkcji r
%% należy do rozdziału 7. To, że cały diagram bifurkacyjny leży wzdłuż linii
%% przez zbiór Mandelbrota -- i przeliczenie r na c -- to wypłata rozdziału 12:
%% ten rozdział wskazuje na nią jednym zdaniem i niczego z niej nie liczy.
%%
%% Bliźniak: chapters/en/ch08-bifurcations.tex. Podrozdział za podrozdziałem,
%% makro za makrem, liczba za liczbą; tools/parity.py je porównuje.

\chapter{Droga do chaosu i liczba, której nikt się nie spodziewał}
\label{ch:bifurcations}
\index{diagram bifurkacyjny}\index{stała Feigenbauma}

\noindent
W rozdziale~\ref{ch:logistic} ustawiłeś jedno pokrętło, $r$, w czterech
pozycjach i zobaczyłeś cztery zachowania. Przy $r = \num{2.8}$ odwzorowanie
logistyczne ustaliło się na jednej wartości; przy \num{3.2} przeskakiwało
między dwiema; przy \num{3.5} krążyło po czterech; przy \num{3.9} nigdy się
nie powtórzyło. Cztery zdjęcia jednej maszyny. Co dzieje się pomiędzy nimi?
Czy porządek zamienia się w chaos przy jednym ustawieniu pokrętła, tak jak
woda wrze w stu stopniach, czy po drodze?

Odpowiedzią jest jeden obraz, jeden z najczęściej reprodukowanych w nauce, a
ukryta jest w nim liczba, której nikt tam nie umieścił. Wychodzi też z innych
równań, zmierzono ją w ciekłym helu, a pod koniec rozdziału sam obliczysz ją
z dokładnością do pięciu cyfr, używając wyłącznie dodawania, odejmowania,
mnożenia i dzielenia.

\begin{outcomes}
\outcome{odczytać diagram bifurkacyjny i powiedzieć, co robi odwzorowanie
  przy każdym pokazanym na nim ustawieniu pokrętła}
\outcome{zlokalizować podwojenie okresu, obserwując przebiegi odwzorowania,
  i powiedzieć, co ogranicza dokładność takiego pomiaru}
\outcome{znaleźć parametr superstabilny metodą bisekcji, używając wyłącznie
  $+$, $-$, $\times$ i $\div$}
\outcome{obliczyć stosunek kolejnych odstępów między podwojeniami i
  oszacować, gdzie podwojenia się kończą}
\outcome{przeprowadzić to samo poszukiwanie dla drugiej reguły i powiedzieć,
  co pokazuje porównanie}
\outcome{znaleźć okno porządku w zakresie chaotycznym i przytoczyć, co Li i
  Yorke udowodnili o okresie trzy}
\end{outcomes}

\section{Jeden obraz dla każdego ustawienia pokrętła}
\label{sec:ch08-picture}
\index{diagram bifurkacyjny}\index{stan przejściowy}

\noindent
Żeby zobaczyć wszystkie $r$ naraz, potrzebujesz przepisu, który zamienia
cały przebieg w kilka kropek, a potem rysujesz kropki tysięcy przebiegów
obok siebie.

\begin{enumerate}
\item Wybierz wartość $r$ i zacznij odwzorowanie od $x = \num{0.5}$.
\item Wykonaj tysiąc kroków i wyrzuć je. Ten odcinek to \emph{stan
  przejściowy}: orbita szuka drogi tam, dokąd zmierza.
\item Wykonaj jeszcze trzysta i nad $r$ narysuj kropkę na wysokości każdej
  wartości.
\end{enumerate}

\pyregion{code/ch08/bifurcation.py}{longrun}{Przepis: pozwól orbicie się ustalić, potem zachowaj to, co odwiedza}{lst:ch08-longrun}

\noindent
\code{orbit} to funkcja z biblioteki książki, która stosuje regułę raz za
razem i zachowuje każdą wartość; \code{xs[burn + 1:]} zostawia te po stanie
przejściowym.

\begin{think}
Przy $r = \num{2.8}$ odwzorowanie ustala się w punkcie stałym $1 - 1/r$,
który znalazłeś w rozdziale~\ref{ch:logistic}. Ile różnych kropek narysuje
przepis nad $r = \num{2.8}$, gdy stan przejściowy zostanie wyrzucony?
\end{think}
\reveal{Jedną. Wszystkie trzysta wartości to \val{ch08.fixed} z dokładnością
do czterech miejsc po przecinku, czyli $1 - 1/\num{2.8}$, więc kropki lądują
jedna na drugiej.}
Cykl dwóch wartości rysuje dwie kropki, cykl czterech \dash{} cztery, a
orbita, która nigdy się nie powtarza, rysuje trzysta różnych: krótką pionową
smugę. Reszta listingu pyta o sześć wartości $r$:

\transcript{ch08-bifurcation}

\noindent
Dwa wiersze są nowe: przy \num{3.56} orbita krąży po ośmiu wartościach, a
przy \num{3.567} po szesnastu. Teraz wykonaj przepis dla tysięcy wartości
$r$ między \num{2.8} a $4$ i narysuj wszystkie kropki naraz.

\plotfig{ch08-diagram}{Diagram bifurkacyjny odwzorowania logistycznego. Nad
każdą wartością $r$ kropki pokazują, gdzie $x$ spędza czas po zakończeniu
stanu przejściowego: jedna linia tam, gdzie się ustala, dwie tam, gdzie
przeskakuje, smuga tam, gdzie nigdy się nie powtarza.}{diagram bifurkacyjny}

\noindent
To jest \emph{diagram bifurkacyjny}, cały rozdział~\ref{ch:logistic} w
jednym kadrze. Czytaj go od lewej do prawej, jak podróż wzdłuż pokrętła.
Jedna linia aż do $r = 3$; tam rozdziela się na dwie; w pobliżu \num{3.45}
każda gałąź rozdziela się znowu; potem wszystko rozpływa się w smugę
poprzecinaną jasnymi pasami. \emph{Bifurkacja} (od łacińskiego słowa
oznaczającego dwuzębne widły) to wartość pokrętła, przy której zmienia się
rodzaj długoterminowego zachowania. Każde rozwidlenie tutaj to
\emph{podwojenie okresu}: każda odwiedzana wartość staje się dwiema, więc
cykl potrzebuje dwa razy więcej czasu, żeby się zamknąć.

\begin{think}
Spójrz na obraz między $r = \num{3.5}$, gdzie orbita ma okres cztery, a
smugą. Czy porządek ustępuje chaosowi przy jednej wartości $r$, jak po
przestawieniu przełącznika?
\end{think}
\reveal{Wygląda jak przełącznik, ale nim nie jest. Powiększ odcinek przed
smugą, a każda gałąź rozdzieli się znowu i znowu: okresy $8$, $16$, $32$ i
tak dalej bez końca, wszystkie mieszczące się przed $r = \val{ch08.end}$.
Tabela złapała już dwa z nich.}

\begin{trapbox}
\textbf{\enquote{Przejście od porządku do chaosu to nagłe przełączenie.}}
Nie na tej drodze. Porządek rozpada się w nieskończonej \emph{kaskadzie}
podwojeń, z których każde przychodzi szybciej niż poprzednie, a chaos
zaczyna się tam, gdzie kaskada się kończy. Smuga wygląda na nagłą tylko
dlatego, że jej ostatnie rozwidlenia leżą zbyt blisko siebie, by je
narysować. Inne układy dochodzą do chaosu innymi drogami; ta,
\emph{kaskada podwojeń okresu}, to droga, którą idzie pojedynczy garb.
\end{trapbox}

\section{Rozwidlenia, które się zagęszczają}
\label{sec:ch08-cascade}
\index{podwojenie okresu}\index{bisekcja}

\noindent
Gdzie dokładnie leży każde rozwidlenie? Tabela wyznacza dla nich
przedziały: okres jeden przy \num{2.8}, okres dwa przy \num{3.2}, więc
pierwsze rozwidlenie leży pomiędzy. Wypróbuj środek przedziału, sprawdź, na
jakim okresie odwzorowanie się tam ustala, zatrzymaj połowę, w której wciąż
jest zmiana, i powtórz. Dwadzieścia rund zmniejsza przedział milion razy. To
jest \emph{bisekcja}, a potrzebuje ona tylko sposobu, żeby zapytać
\enquote{czy tutaj to wciąż okres $p$?}

\pyregion{code/ch08/doublings.py}{fork}{Namierzanie rozwidlenia przez połowienie przedziału}{lst:ch08-fork}

\noindent
\code{settled\_period} wykonuje dwadzieścia tysięcy kroków odwzorowania i
pyta biblioteczną funkcję \code{period}, co się powtarza. Listing uruchamia
\code{fork} na czterech przedziałach, które daje tabela.

\begin{think}
Okazuje się, że pierwsze dwa rozwidlenia dzieli \val{ch08.gap.one} w $r$, a
kolejne dwa \val{ch08.gap.two}. Ile razy mniejszy jest drugi odstęp? Jeśli
następny zmaleje o ten sam czynnik, ile będzie wynosił?
\end{think}
\reveal{Około \val{ch08.ratio.one} razy mniejszy:
$\val{ch08.gap.one} \div \val{ch08.gap.two}$. Ten sam czynnik jeszcze raz
dałby następny odstęp około \num{0.02}.}
Oto, co zmierzył listing:

\transcript{ch08-doublings}

\noindent
Trzeci odstęp wynosi \val{ch08.gap.three}; drugi podzielony przez trzeci
daje \val{ch08.ratio.two}. Każdy odstęp jest więc nieco mniej niż pięć razy
większy od następnego, a pozostałe odstępy (ostatni podzielony przez
\num{4.7}, potem jeszcze raz przez \num{4.7} i tak dalej) dają razem
ostatni odstęp podzielony przez $\num{4.7} - 1$, czyli nieco ponad jego
ćwierć. Nieskończenie wiele rozwidleń mieści się na skończonym odcinku
pokrętła: kaskada kończy się niedaleko za \val{ch08.fork.four}.

\begin{think}
Pierwsze rozwidlenie znamy dokładnie. W rozdziale~\ref{ch:logistic} okazało
się, że punkt stały traci stabilność przy $r = 3$, gdzie jego nachylenie,
$2 - r$, przekracza $-1$. Tabela podaje \val{ch08.fork.one}. Co poszło nie
tak?
\end{think}
\reveal{Nic w regule: pomiarowi zabrakło cierpliwości. Przy
$r = \num{2.9998}$ każdy krok mnoży małe zaburzenie przez
$2 - r = \num{-0.9998}$, więc zaburzenie maleje o jedną pięćdziesiątą
procenta na krok, a po dwudziestu tysiącach kroków orbita wciąż wyraźnie
przeskakuje z boku na bok. \code{settled\_period} uznaje to za cykl o
okresie dwa.}
Każde rozwidlenie mierzone przez obserwację wypada za wcześnie z tego
samego powodu: blisko rozwidlenia cykl, który traci stabilność, ustala się
bardzo powoli.

\begin{note}
Drugie rozwidlenie da się znaleźć dokładnie. Przy jednym obiegu cyklu o
okresie dwa zaburzenie jest mnożone przez nachylenie w każdym z jego dwóch
punktów, a strona szkolnej algebry pokazuje, że iloczyn wynosi
$4 + 2r - r^{2}$. Osiąga on $-1$, gdy $r^{2} - 2r - 5 = 0$, czyli przy
$r = 1 + \sqrt{6} = \val{ch08.fork.two.exact}$; \val{ch08.fork.two} z tabeli
wypada za wcześnie, zgodnie z oczekiwaniem. Późniejsze rozwidlenia nie mają
tak prostego wzoru.
\end{note}

\section{Szczyt garbu}
\label{sec:ch08-superstable}
\index{cykl superstabilny}\index{Feigenbaum, Mitchell}

\noindent
W 1975 roku Mitchell Feigenbaum, fizyk z laboratorium w Los Alamos w Nowym
Meksyku, wyznaczał położenie tych rozwidleń na programowalnym kalkulatorze
kieszonkowym i zauważył wzór, który właśnie zobaczyłeś: każdy odstęp to
stała część poprzedniego. Żeby zmierzyć tę część dokładnie, użyj lepiej
położonego punktu orientacyjnego niż rozwidlenie, na którym cykl trzyma się
resztkami sił: punktu na każdym odcinku, w którym cykl trzyma się
najmocniej.

Nachylenie odwzorowania logistycznego w punkcie $x$ wynosi $r(1 - 2x)$
\dash{} to współczynnik rozciągnięcia, który nauczyłeś się mierzyć w
rozdziale~\ref{ch:feedback}. Jest równe zeru przy $x = \num{0.5}$, na
szczycie garbu. Cykl przechodzący przez $\num{0.5}$ mnoży każde zaburzenie
przez to zero raz na obieg i je wymazuje. Taki cykl nazywa się
\emph{superstabilnym}, a dla każdego okresu $1, 2, 4, 8, \ldots$ jest na
drodze dokładnie jedna wartość $r$, przy której to się dzieje: w tym
rozdziale nazywamy te wartości \emph{punktami orientacyjnymi}. Każdy z nich
to rozwiązanie jednego równania: zacznij od $\num{0.5}$, zrób tyle kroków,
ile wynosi okres, i wróć do $\num{0.5}$.

\begin{think}
Dla okresu jeden: przy jakim $r$ jeden krok przenosi $\num{0.5}$ z powrotem
prosto do $\num{0.5}$? Krok to $r \times \num{0.5} \times \num{0.5}$.
\end{think}
\reveal{$r = 2$. Połowa połowy to ćwierć, a $2 \times \num{0.25} =
\num{0.5}$.}
Dla dłuższych okresów nie ma skrótu w pamięci, a bisekcja żadnego nie
potrzebuje. \code{from\_top} podaje, jak daleko od $\num{0.5}$ kończy się
orbita; ta odległość jest dodatnia po jednej stronie punktu orientacyjnego i
ujemna po drugiej:

\pyregion{code/ch08/superstable.py}{search}{Zacznij na szczycie i połowiaj przedział, aż orbita do niego wróci}{lst:ch08-search}

\noindent
Pierwsze dwa punkty orientacyjne otaczamy przedziałami na oko, z obrazu.
Potem każdy odstęp jest mniejszy od poprzedniego, więc następnego punktu
szukamy między jedną dziesiątą a połową ostatniego odstępu dalej \dash{} to
hojny przedział dla odstępów malejących o czynnik bliski pięciu.
\emph{Rodzina} to funkcja, która dla danego $r$ buduje regułę.

\pyregion{code/ch08/superstable.py}{landmarks}{Punkty orientacyjne dla okresów $1, 2, 4, 8, \ldots$, każdy otoczony na podstawie poprzedniego}{lst:ch08-landmarks}

\noindent
Listing znajduje pierwszych \val{ch08.landmark.count} punktów
orientacyjnych, dla okresów od $1$ do $1024$, używając wyłącznie $+$, $-$,
$\times$ i $\div$, więc każda cyfra jest taka sama na każdym komputerze:

\transcript{ch08-superstable}

\noindent
Spójrz w dół ostatniej kolumny: każdy odstęp podzielony przez następny.
Przez kilka wierszy się waha, a potem przestaje się zmieniać:
\val{ch08.delta}. To \emph{stała Feigenbauma}, oznaczana $\delta$ (grecka
litera delta). Punkty orientacyjne dają ją znacznie dokładniej niż
rozwidlenia, bo w punkcie orientacyjnym orbita ustala się od razu, a nie
ledwie wcale. Rozwidlenia i punkty orientacyjne występują na drodze na
zmianę, po jednym z każdego na okres, więc jedne i drugie zagęszczają się w
tym samym tempie ku temu samemu końcowi, $r = \val{ch08.end}$. Punkt
orientacyjny okresu dwa, \val{ch08.super.two}, to dokładnie $1 + \sqrt{5}$,
gdzie iloczyn z notatki, $4 + 2r - r^{2}$, jest równy zeru.

Feigenbaum opublikował wyniki w 1978 roku; Pierre Coullet i Charles Tresser
we Francji znaleźli tę samą liczbę niezależnie, mniej więcej w tym samym
czasie. To, co zrobił potem, sprawiło, że ta liczba jest sławna.

\begin{lab}{l08_01_superstable}{Znajdź punkt orientacyjny sam}
Napisz \code{from\_top} i \code{bisect}, a potem połącz je w
\code{superstable}. Sprawdź swoją odpowiedź dla okresu dwa, porównując ją z
$1 + \sqrt{5}$. Następnie znajdź $r$ między \num{3.82} a \num{3.84}, przy
którym $\num{0.5}$ wraca do siebie po \emph{trzech} krokach. Jeszcze
spotkasz tę liczbę.
\end{lab}

\section{Liczba, która należy do drogi}
\label{sec:ch08-universal}
\index{uniwersalność}\index{odwzorowanie sinusowe}

\noindent
Następnym krokiem Feigenbauma było wypróbowanie innej reguły,
\emph{odwzorowania sinusowego}
\[ x_{n+1} = s \sin(\pi x_{n}), \]
które zwraca $s$ razy sinus z $\pi x$. Nie potrzebujesz trygonometrii, żeby
za nim nadążyć: dla $x$ między $0$ a $1$ sinus z $\pi x$ to gładki garb,
który rośnie od $0$ do $1$ przy $x = \num{0.5}$ i opada z powrotem do $0$.
Odwzorowanie sinusowe ma więc kształt odwzorowania logistycznego \dash{}
jeden garb ze szczytem w $\num{0.5}$ \dash{} i nic z jego wzoru: żadnego
$x(1 - x)$, inne pokrętło $s$, krzywą, która nie jest parabolą.

\begin{think}
Przeprowadź to samo poszukiwanie punktów orientacyjnych dla odwzorowania
sinusowego. Czy spodziewasz się, że jego odstępy będą maleć znowu o
\val{ch08.delta} z odwzorowania logistycznego, czy o jakąś inną liczbę,
należącą do sinusa?
\end{think}
\reveal{O tę samą liczbę. Odstępy odwzorowania sinusowego maleją najpierw o
\val{ch08.sine.first}, a potem, punkt po punkcie, zbliżają się do
\val{ch08.sine.delta}, na której zatrzymują się też stosunki odwzorowania
logistycznego.}
Listing nie wie o odwzorowaniu sinusowym nic poza jego rodziną i dwoma
przedziałami odczytanymi z jego diagramu:

\pyregion{code/ch08/two_maps.py}{compare}{To samo poszukiwanie dla dwóch niepowiązanych reguł}{lst:ch08-compare}

\transcript{ch08-two-maps}

\noindent
Obie drogi kończą się w różnych miejscach (ostatni wiersz: \val{ch08.end} i
\val{ch08.sine.end}) i inaczej się zaczynają, a jednak zatrzymują się na tym
samym stosunku, co do każdej wydrukowanej cyfry. Funkcja \code{ratios} dzieli
każdy odstęp przez następny; drugie laboratorium prosi cię, żebyś ją
napisał.

\plotfig{ch08-ratios}{Każdy odstęp między punktami orientacyjnymi podzielony
przez następny, dla dwóch reguł, które łączy tylko pojedynczy gładki garb.
Zaczynają daleko od siebie i kończą na tej samej liczbie.}{stosunki dla
dwóch reguł zbiegają się}

\begin{trapbox}
\textbf{\enquote{Inne równanie daje inną drogę.}} Daje inne drogowskazy
\dash{} rozwidlenia leżą przy innych wartościach pokrętła \dash{} ale droga
ma ten sam kształt, a jej rozwidlenia zagęszczają się w tym samym tempie. Ta
liczba należy do drogi, nie do równania. To jest \emph{uniwersalność} i
dlatego stała przyrody może wyjść z zabawkowego modelu.
\end{trapbox}

\noindent
Dlaczego? O każdym kolejnym rozwidleniu decyduje coraz krótszy odcinek
odwzorowania w pobliżu szczytu garbu, oglądany przez coraz silniejsze szkło
powiększające. Powiększ dowolny gładki, zaokrąglony garb dostatecznie mocno
wokół szczytu, a wygląda jak parabola; wszystko, co czyniło odwzorowanie
sinusowe sinusem, zostało powiększone poza pole widzenia. Stałe szkolnej
geometrii, takie jak $\pi$, należą do kształtów. Ta należy do procesu.

\mermaidfig{ch08-universal}{Uniwersalność w jednym zdaniu: każda reguła z
jednym gładkim garbem idzie tą samą drogą, a droga niesie własny
stosunek.}{każdy garb, jedna droga, jeden stosunek}

\noindent
Był to szok, bo przewidywał coś o układach, których równań nikt nie potrafił
zapisać: jeśli któryś dochodzi do chaosu przez podwajanie okresu, jego
rozwidlenia powinny się zagęszczać w tempie Feigenbauma, bez względu na to, z
czego jest zrobiony.

\begin{worldbox}
Około 1980 roku Albert Libchaber i Jean Maurer w Paryżu podgrzewali od dołu
maleńkie pudełko z ciekłym helem i rejestrowali temperaturę w jednym
punkcie, gdy płyn się kotłował. Gdy zwiększali grzanie, rytm temperatury
podwoił okres, a potem podwoił go znowu. Późniejsze doświadczenia z
konwekcją w cieczach i z obwodami elektronicznymi mierzyły stosunek
kolejnych odstępów i znajdowały liczby bliskie liczbie Feigenbauma: nie z
dokładnością do czterech miejsc po przecinku, bo w laboratorium da się
rozróżnić tylko kilka pierwszych podwojeń, ale na tyle bliskie, że nikt nie
uważa tego za przypadek. W 1986 roku Feigenbaum i Libchaber otrzymali
wspólnie Nagrodę Wolfa w dziedzinie fizyki.
\end{worldbox}

\begin{rigourbox}
Nie dowodzimy tutaj, dlaczego każdy gładki garb daje tę samą stałą
$\delta$. Feigenbaum wyjaśnił to metodą z fizyki, \emph{renormalizacją},
która porównuje odwzorowanie po każdym podwojeniu z przeskalowaną kopią
samego siebie; Oscar Lanford podał dowód wspomagany komputerowo w 1982 roku.
\emph{Gładkość} ma znaczenie: garb o szczycie bardziej płaskim niż u
paraboli daje inną stałą.
\end{rigourbox}

\begin{lab}{l08_02_ratio}{Stosunek Feigenbauma i miejsce, w którym kończą się podwojenia}
Napisz \code{ratios}, które dzieli każdy odstęp przez następny, oraz
\code{where\_it\_ends}, które sumuje wszystkie pozostałe odstępy, zakładając,
że każdy to poprzedni podzielony przez ostatni stosunek. Test podaje twoim
funkcjom punkty orientacyjne obu odwzorowań: jeden stosunek, dwa różne
końce.
\end{lab}

\section{Porządek wewnątrz chaosu}
\label{sec:ch08-windows}
\index{okno okresowe}\index{twierdzenie Li--Yorke'a}

\noindent
Za $r = \val{ch08.end}$ podwojenia się wyczerpują, a diagram zamienia się w
smugę.

\begin{think}
Czy od $r = \val{ch08.end}$ do $r = 4$ odwzorowanie logistyczne jest
chaotyczne przy każdym ustawieniu pokrętła, tak jak przy \num{3.9}?
\end{think}
\reveal{Nie. Spójrz jeszcze raz na obraz: smugę przecinają jasne pasy, a w
każdym z nich orbita wraca do krótkiego cyklu. Najszerszy, w pobliżu
\num{3.83}, to okno okresu trzy.}
\code{code/ch08/windows.py} przeprowadza $r$ przez ten pas z tą samą funkcją
\code{settled\_period}, a potem szuka punktów orientacyjnych w jego
wnętrzu:

\transcript{ch08-windows}

\noindent
Przy \num{3.8284} orbita się nie powtarza; przy \num{3.8285} ustaliła się
już w cyklu o okresie trzy. Okno otwiera się pomiędzy nimi, dokładnie przy
$r = 1 + \sqrt{8} = \val{ch08.window.open}$. Cykl o okresie trzy podwaja się
potem do sześciu i dwunastu, a przy \num{3.85} znów rozpływa się w smugę.
Cała droga powtarza się w miniaturze. Pierwszy punkt orientacyjny okna to
$r$, które znalazłeś w pierwszym laboratorium, \val{ch08.super.three}, a jego
punkty orientacyjne zagęszczają się w stosunku \val{ch08.window.delta}:
znowu liczba Feigenbauma, tym razem dla okresów $3, 6, 12, \ldots$ zamiast
$1, 2, 4, \ldots$.

\plotfig{ch08-zoom}{Po lewej: okno okresu trzy. Po prawej: środkowa gałąź
okna w powiększeniu: mała kopia całego diagramu bifurkacyjnego, z własną
kaskadą, własną smugą i własnymi oknami.}{okno okresu trzy w powiększeniu}

\begin{trapbox}
\textbf{\enquote{Gdy układ raz osiągnie chaos, jest chaotyczny wszędzie
dalej.}} Między $r = \val{ch08.end}$ a $4$ odwzorowanie logistyczne wraca do
porządku nieskończenie wiele razy, w oknach o każdym okresie. Przekręć
pokrętło odrobinę, a chaotyczny układ może stać się idealnie okresowy;
przekręć je odrobinę dalej, a znów jest chaotyczny. Wykładnik z
rozdziału~\ref{ch:lyapunov} spada poniżej zera dokładnie w tych pasach.
\end{trapbox}

\begin{note}
W 1975 roku Tien-Yien Li i James Yorke opublikowali pracę, której tytuł był
jej wynikiem: \emph{Period three implies chaos} (okres trzy pociąga za sobą
chaos). W skrócie: jeśli ciągła reguła przenosi przedział w siebie, tak jak
odwzorowanie logistyczne przenosi przedział od $0$ do $1$, i jakiś punkt
startowy wraca do siebie po dokładnie trzech krokach, to reguła ma cykle
każdej długości oraz punkty startowe, których orbity nigdy nie ustalają się
w żadnym cyklu. Praca nadała słowu \emph{chaos} jego matematyczne znaczenie;
silniejszy wynik udowodnił w 1964 roku Ołeksandr Szarkowski w Kijowie, ale
na Zachodzie był mało znany aż do czasu tej pracy. W oknie prawie każdy
punkt startowy kończy na cyklu o okresie trzy, ale twierdzenie nadal
obowiązuje: pozostałe cykle są na miejscu, niestabilne, a orbita może długo
błąkać się między nimi, zanim się ustali.
\end{note}

\noindent
Diagram ma jeszcze jedną niespodziankę, a należy ona do
rozdziału~\ref{ch:mandelbrot}: cały ten obraz leży wzdłuż linii
przecinającej najsłynniejszy obraz matematyki, a rozwidlenia i okno są
dokładnie tam, gdzie umieszcza je jednolinijkowe przeliczenie.

\begin{summarybox}
\begin{itemize}
\item \textbf{Diagram bifurkacyjny} pokazuje dla każdej wartości pokrętła,
  gdzie orbita przebywa po zakończeniu stanu przejściowego.
  \textbf{Bifurkacja} to wartość, przy której zmienia się rodzaj
  długoterminowego zachowania; na tej drodze każda jest \textbf{podwojeniem
  okresu}.
\item Rozwidlenia się zagęszczają, każdy odstęp jest nieco mniej niż pięć
  razy większy od następnego, więc nieskończenie wiele z nich mieści się
  przed $r = \val{ch08.end}$. Zmierzone przez obserwację rozwidlenie wypada
  za wcześnie, bo w jego pobliżu orbita ustala się bardzo powoli.
\item W \textbf{superstabilnym} punkcie orientacyjnym cykl przechodzi przez
  szczyt garbu, gdzie nachylenie jest zerowe. Każdy taki punkt to jedno
  równanie, rozwiązywane \textbf{bisekcją} za pomocą $+$, $-$, $\times$ i
  $\div$, a stosunek odstępów między nimi zatrzymuje się na \textbf{stałej
  Feigenbauma}, \val{ch08.delta}.
\item Odwzorowanie sinusowe daje ten sam stosunek: liczba należy do drogi,
  którą idzie każda gładka reguła z jednym garbem, a nie do wzoru. To jest
  \textbf{uniwersalność}, zmierzona w cieczach i obwodach.
\item Za kaskadą chaos przerywają \textbf{okna} porządku. Najszersze, okresu
  trzy, otwiera się przy $1 + \sqrt{8}$ i zawiera miniaturę całej drogi; Li
  i Yorke udowodnili, że okres trzy pociąga za sobą chaos.
\end{itemize}
\end{summarybox}

\canyou

\begin{checkpoint}
\item Co pokazuje kolumna kropek nad jedną wartością $r$ na diagramie
  bifurkacyjnym i jak wygląda dla punktu stałego, cyklu o okresie dwa i
  chaosu?
  \answerto{Gdzie orbita przebywa na dłuższą metę przy tym $r$, po stanie
  przejściowym: jedna kropka dla punktu stałego, dwie dla cyklu o okresie
  dwa, smuga dla chaosu.}
\item Dwa kolejne odstępy między rozwidleniami wynoszą \num{0.08} i
  \num{0.017}. Mniej więcej jaki jest ich stosunek i ile powinien wynosić
  następny odstęp?
  \answerto{Około \num{4.7}, z $\num{0.08} \div \num{0.017}$; następny odstęp
  powinien więc wynosić około $\num{0.017} \div \num{4.7}$, czyli mniej
  więcej \num{0.0036}.}
\item Dlaczego rozwidlenie zlokalizowane przez obserwację orbity wypada
  nieco za wcześnie?
  \answerto{Blisko rozwidlenia cykl, który traci stabilność, mnoży
  zaburzenie przy każdym obiegu przez prawie $-1$, więc ustala się niezwykle
  powoli, a przebieg o skończonej długości wciąż pokazuje przeskakiwanie
  między wartościami.}
\item Czym jest superstabilny punkt orientacyjny i dlaczego da się go
  znaleźć z dokładnością do wielu cyfr za pomocą samych $+$, $-$, $\times$ i
  $\div$?
  \answerto{To $r$, przy którym cykl przechodzi przez szczyt garbu,
  $x = \num{0.5}$, gdzie nachylenie jest zerowe. Jest rozwiązaniem jednego
  równania (zacznij od $\num{0.5}$, zrób tyle kroków, ile wynosi okres, wróć
  do $\num{0.5}$), które bisekcja rozwiązuje, połowiąc przedział.}
\item Podwojenia odwzorowania sinusowego kończą się około
  \val{ch08.sine.end}, a logistycznego około \val{ch08.end}. Czy to przeczy
  uniwersalności?
  \answerto{Nie. Uniwersalność mówi, że stosunek kolejnych odstępów dąży do
  tej samej liczby, około \val{ch08.delta}, dla każdej reguły z jednym
  gładkim garbem; to, gdzie kończy się każda droga, zależy od reguły.}
\item Czy odwzorowanie logistyczne jest chaotyczne przy $r = \num{3.83}$? Co
  Li i Yorke udowodnili o okresie trzy?
  \answerto{Nie: \num{3.83} leży w oknie okresu trzy, które otwiera się przy
  $1 + \sqrt{8}$, około \val{ch08.window.open}. Li i Yorke udowodnili, że
  ciągła reguła przenosząca przedział w siebie, mająca punkt o okresie trzy,
  ma cykle każdej długości i punkty startowe, które nigdy nie ustalają się w
  żadnym cyklu.}
\end{checkpoint}
